%Präambel
\documentclass[fontsize=12, a4aper]{scrartcl}
\usepackage[english, ngerman]{babel}
\usepackage{mwe}
\usepackage{newfloat}
\usepackage[utf8]{inputenc}
\usepackage[autostyle=true,german=quotes]{csquotes}
\usepackage[T1]{fontenc}
\usepackage{amssymb}

\usepackage{amsmath}

\usepackage[font=footnotesize,labelfont=bf,singlelinecheck=false]{caption}

\usepackage{tabularx}
\usepackage{subfig}
\usepackage{caption, booktabs}

\usepackage{float}

\DeclareFloatingEnvironment
[
fileext=loa,
listname=Anlagenverzeichnis,
name=Anlage,
placement=htbp,
within=section,
]{anlage}

\newcommand\textbox[1]
{
	\parbox{.333\textwidth}{#1}
}

\newcommand\Tstrut{\rule{0pt}{3ex}}           % = `top' strut
\newcommand\Bstrut{\rule[-0.9ex]{0pt}{0pt}}   % = `bottom' strut

\begin{document}
	
	\begin{table}[H]
		
		\fontsize{10pt}{12pt}
		
		\caption{Testfunktionen zum Vergleich mit Algorithmen, die probabilistisches Lernen einsetzen}
		\label{einfacheTestfunktionen}
		
		\begin{tabularx}{\textwidth}{l|r}
			
			\midrule
			
			Funktion, Definitionsbereich, & Formel\\
			Optimum & \\
			
			\midrule
			
			Cigar, \overrightarrow{x} \(\in [-3,7]^n\), n = 10, & \(f(x)=x_1^2 + 10^4\sum_{i=2}^{n} x_i^2 )\)\Tstrut\\
			0 bei \overrightarrow{x} = [0, 0, ...] & \\
			
			Diagonal plane, \overrightarrow{x} \(\in [0.5,1.5]^n\), n = 10, & \(\frac{1}{n} \sum_{i=1}^{n} x_i\)\Tstrut\\
			\(10^{10}\) bei mehreren & \\
			
			Ellipsoid, \overrightarrow{x} \(\in [-3,7]^n\), n = 10, & \(f(x)=\sum_{i=1}^{n}(100^\frac{i-1}{n-1} x_i)\)\Tstrut\\
			0 bei \overrightarrow{x} = [0, 0, ...] & \\
			
			Paraboloid, \overrightarrow{x} \(\in [-3,7]^n\), n = 10, & \(f(x)=\sum_{i=1}^{n} x^2\)\Tstrut\\
			0 bei \overrightarrow{x} = [0, 0, ...] & \\
			
			Plane, \overrightarrow{x} \(\in [0.5,1.5]^n\), n = 10, & \(f(x)=x_1\)\Tstrut\\
			\(10^{10}\) bei \(x_1 = 10^{10}\) & \\
			
			Rosenbrock, \overrightarrow{x} \(\in [-3,7]^n\), n = 10,  & \(f(x)=\sum_{i=1}^{n-1} 100 (x^2_i-x_{i+1})^2  + (x_i-1)^2\) \Tstrut\\
			0 bei \overrightarrow{x} = [1, 1, ...] & \\
			
			Tablet, \overrightarrow{x} \(\in [-3,7]^n\), n = 10 & \(f(x)=10^4x_1^2 + \sum_{i=2}^{n}x^2_i \)\Tstrut\\
			0 bei \overrightarrow{x} = [0, 0, ...] & 
			
		\end{tabularx}
		
	\end{table}
	
	\bigskip
	
	\begin{table}[H]
		
		\fontsize{10pt}{12pt}
		
		\caption{erster Teil der Testfunktionen zum Vergleich mit anderen Ameisenalgorithmen verwandten Methoden und Metaheuristiken, die für stetige Probleme adaptiert wurden}
		\label{mittlereTestfunktionen}
		
		\begin{tabularx}{\textwidth}{l|r}
			
			\midrule
			
			Funktion, Definitionsbereich, & Formel\\
			Optimum & \\
			
			\midrule
			
			\textit{\(B_2\)}, \overrightarrow{x} \(\in [-100,100]^n\), n = 2 & \(f(x)=x_1^2 + 2x_2^2 - \frac{3}{10} cos(3 \pi x_1) -\frac{2}{5} cos(4 \pi x_2) +7/10\)\\
			0 bei \overrightarrow{x} = [0, 0] &\\
			
			Branin RCOS, \overrightarrow{x} \(\in [-5,15]^n\), n = 2
			& \(f(x)=(x_2 - \frac{5}{4\pi}^2 x_1^2 + \frac{5}{\pi} x_1 - 6)^2 + 10 (1 - \frac{1}{8\pi})cos(x_1) + 10\)\Tstrut\\
			0.397887 bei \overrightarrow{x} = [-\(\pi\) , 12.275] & \\ 
			od. x = [\(\pi\) , 2.275] & \\
			od. x = [9.42478, 2.475] & \\
			
			Easom, \overrightarrow{x} \(\in [-100,100]^n\), n = 2 & \(f(x)=-cos(x_1) \, cos(x_2) e^{(x_1-\pi)^2(x_2-\pi)^2}\)\Tstrut\\
			0 bei \overrightarrow{x} = \([\pi \cdot k + \pi/2\), & \\
			\( \pi \cdot k + \pi/2], k \in \mathbb{Z}\) & \\
			
			De Jong, \overrightarrow{x} \(\in [-5.12,5.12]^n\), n = 3,
			& \(f(x)=x_1^2 + x_2^2 + x_3^2\)\Tstrut\\
			0 bei \overrightarrow{x} = [0, 0, 0] &\\
			
			Goldstein, \overrightarrow{x} \(\in [-2,2]^n\) & \(f(x)=(1 + (x_1 + x_2 + 1)^2 (19 - 14x_1 + 13x_1^2\)\Tstrut\\
			and Price, n = 2 & \( - 14x_2 + 6x_1x_2 + x_2^2))(30 + (2x_1 - 3x_2)^2 \cdot\)\\
			3 bei \overrightarrow{x} = [0,1] & \((18 - 32x_1 + 12x_1^2 - 48x_2 - 36x_1x_2 - 27x_2^2))\)\\
			
			Griewangk, \overrightarrow{x} \(\in [-600, 600]^n, n = 10\) & \(f(x) = \sum_{i=1}^{n} \frac{x_i^2}{4000} - \prod_{i=1}^{n} \cos(\frac{x_i}{\sqrt{i}}) - 1\) \Tstrut\\
			0 bei \overrightarrow{x} = [0, 0, ...] & 
			
		\end{tabularx}
		
	\end{table}
	
	\bigskip
	
	\begin{table}[H]
		
		\fontsize{10pt}{12pt}
		
		\caption{zweiter Teil der Testfunktionen zum Vergleich mit anderen Ameisenalgorithmen verwandten Methoden und Metaheuristiken, die für stetige Probleme adaptiert wurden}
		\label{schwereTestfunktionen}
		
		\begin{tabularx}{\textwidth}{l|r}
			
			\midrule
			Funktion, Definitionsbereich, & Formel\\
			Optimum & \\
			
			\midrule
			
			Hartmann(3,4), \overrightarrow{x} \(\in [0,1]^n\), & \(f_{H_{3,4}}(x) = - \sum_{i = 1}^{4} c_i e^{-\sum_{j=1}^{3} a_{ij} (x_j - p_{ij})^2} \)\Tstrut\\
			
			n = 3 & \(a_{ij}\) = 	\(\begin{pmatrix}
			3.0 & 10.0 & 30.0\\
			0.1 & 10.0 & 35.0\\
			3.0 & 10.0 & 30.0\\
			0.1 & 10.0 & 35.0
			\end{pmatrix}\)
			
			\(c_i\) =			\(\begin{pmatrix}
			1.0\\
			1.2\\
			3.0\\
			3.2
			\end{pmatrix}\)\\
			
			& \(p\) = 		\(\begin{pmatrix}
			0.3689 & 0.1170 & 0.2673\\
			0.4699 & 0.4387 & 0.7470\\
			0.1091 & 0.8732 & 0.5547\\
			0.0381 & 0.5743 & 0.8828
			\end{pmatrix}\)\\
			
			Hartmann(6,4), \overrightarrow{x} \(\in [0,1]^n\), & \(f_{H_{3,4}}(x) = - \sum_{i = 1}^{4} c_i e^{-\sum_{j=1}^{6} a_{ij} (x_j - p_{ij})^2} \)\Tstrut\\
			
			n = 6 & \(a_{ij}\) = 	\(\begin{pmatrix}
			10.0 & 3.0 & 17.0 & 3.5 & 1.5 & 8.0\\
			0.05 & 10.0 & 17.0 & 0.1 & 8.0 & 14.0\\
			3.0 & 3.5 & 1.7 & 10.0 & 17.0 & 8.0\\
			17.0 & 8.0 & 0.05 & 10.0 & 0.1 & 14.0
			\end{pmatrix}\)
			
			\(c_i\) =		\(\begin{pmatrix}
			1.0\\
			1.2\\
			3.0\\
			3.2
			\end{pmatrix}\)\\
			
			& \(p\) = 		\(\begin{pmatrix}
			0.1312 & 0.1696 & 0.5569 & 0.0124 & 0.8283 & 0.5886\\
			0.2329 & 0.4135 & 0.8307 & 0.3736 & 0.1004 & 0.9991\\
			0.2348 & 0.1451 & 0.3522 & 0.2883 & 0.3047 & 0.6650\\
			0.4047 & 0.8828 & 0.8732 & 0.5743 & 0.1091 & 0.0381
			\end{pmatrix}\)\\
			Shekel(4,k; k = 5,7,10) & \(f(x) = -\sum_{i=1}^{k}(\sum_{j=1}^{4} (x_j - A_{ji})^2 + c_i)^{-1}\)\Tstrut\\
			\overrightarrow{x}, \(\in [0,1]^n\), n = 6 &
			
			\(A\) = 		\(\begin{pmatrix}
			4.0 & 4.0 & 4.0 & 4.0\\
			1.0 & 1.0 & 1.0 & 1.0\\
			8.0 & 8.0 & 8.0 & 8.0\\
			6.0 & 6.0 & 6.0 & 6.0\\
			3.0 & 7.0 & 3.0 & 7.0\\
			2.0 & 9.0 & 2.0 & 9.0\\
			5.0 & 5.0 & 3.0 & 3.0\\
			8.0 & 1.0 & 8.0 & 1.0\\
			6.0 & 2.0 & 6.0 & 2.0\\
			7.0 & 3.6 & 7.0 & 3.6
			\end{pmatrix}^T\)
			
			\(c_i\) =		\(\begin{pmatrix}
			0.1\\
			0.2\\
			0.2\\
			0.4\\
			0.4\\
			0.6\\
			0.3\\
			0.7\\
			0.5\\
			0.5
			\end{pmatrix}\)
			
		\end{tabularx}
		
	\end{table}
	
\begin{table}[H]
	
	\fontsize{10pt}{12pt}
	
	\caption{zweiter Teil der Testfunktionen zum Vergleich mit anderen Ameisenalgorithmen verwandten Methoden und Metaheuristiken, die für stetige Probleme adaptiert wurden}
	\label{schwereTestfunktionen}
	
	\begin{tabularx}{\textwidth}{r|l}
		
		\fontsize{10pt}{12pt}
		
		
		
		$\begin{pmatrix}
		1 & 2
		\end{pmatrix}$ & \\
		
		\(\begin{pmatrix}
		1 & 2
		\end{pmatrix}\)
		
	\end{tabularx}

\end{table}

\cite{kern2004learning}

\bibliography{literatur.bib}

\bibliographystyle{plain}

\end{document}